% Base editorial: turabian-researchpaper. Las adaptaciones del SRL
% se separan de las opciones que el autor cambia por ensayo.
\usepackage{iftex}
\RequireXeTeX
\usepackage{fontspec}
\usepackage[spanish,es-noshorthands,es-nodecimaldot]{babel}
\usepackage{csquotes}
\DeclareQuoteStyle{spanish}{«}{»}{“}{”}
\setquotestyle{spanish}
\usepackage{xcolor}
\usepackage{graphicx}
\usepackage{etoolbox}
\usepackage{ragged2e}
\usepackage{geometry}
\usepackage{setspace}
\usepackage{titlesec}
\usepackage{titletoc}
\usepackage{fancyhdr}
\usepackage{refcount}
% No se simulan negritas ni cursivas: las cuatro variantes deben existir.
\defaultfontfeatures{Ligatures=TeX}
\newcommand{\SRLLoadSystemFont}{%
  \IfFontExistsTF{\SRLFontName}{%
    \setmainfont{\SRLFontName}%
    \edef\SRLActualFont{\SRLFontName}%
    \typeout{SRL-FONT-SOURCE: system}%
  }{%
    \setmainfont{\SRLFallbackFont}%
    \edef\SRLActualFont{\SRLFallbackFont}%
    \PackageWarningNoLine{srl-template}{Fuente solicitada '\SRLFontName' no disponible.
      Se usa '\SRLFallbackFont'. Esto NO es una compilacion con Minion Pro}%
    \typeout{SRL-FONT-SOURCE: fallback}%
  }%
}
\newcommand{\SRLLoadLocalFont}{%
  \setmainfont{\SRLFontRegular}[
    Path=\SRLFontDirectory,
    UprightFont=\SRLFontRegular,
    ItalicFont=\SRLFontItalic,
    BoldFont=\SRLFontBold,
    BoldItalicFont=\SRLFontBoldItalic
  ]%
  \edef\SRLActualFont{\SRLFontName}%
  \typeout{SRL-FONT-SOURCE: files}%
}
\ifdefstrequal{\SRLFontName}{\SRLLocalFontFamily}{%
  \IfFileExists{\SRLFontDirectory\SRLFontRegular}
    {\SRLLoadLocalFont}{\SRLLoadSystemFont}%
}{\SRLLoadSystemFont}
\AtBeginDocument{\typeout{SRL-FONT: \SRLActualFont}}

% Valores proporcionados por el autor; conservar nombres y RGB exactos.
\definecolor{redSRL}{RGB}{177,35,51}
\definecolor{yellowSRL}{RGB}{252,186,99}
\definecolor{darkgreySRL}{RGB}{237,235,232}
\definecolor{lightgreySRL}{RGB}{247,245,242}

% Curso y fecha permanecen negros en ambos modos.
\colorlet{SRLCoverDetailColor}{black}
\ifSRLColors
  \colorlet{SRLHeadingColor}{redSRL}
  \typeout{SRL-COLORS: true}
\else
  \colorlet{SRLHeadingColor}{black}
  \typeout{SRL-COLORS: false}
\fi


% La guía SRL exige NB; no se sustituye por chicago.bst (autor-fecha).
\usepackage[
  notes,backend=biber,autocite=footnote,
  autolang=none,
  noibid,short=false,sorting=nty,sortlocale=es_ES,
  maxcitenames=3,mincitenames=1,maxbibnames=99,minbibnames=99,
  isbn=false
]{biblatex-chicago}
\DeclareLanguageMapping{spanish}{cms-spanish}
\usepackage{hyperref}
\hypersetup{
  hidelinks,unicode=true,
  pdftitle={\EssayTitle},pdfauthor={\EssayAuthor},
  pdfsubject={\EssayCourse},pdfcreator={XeLaTeX / plantilla SRL}
}
\urlstyle{same}
% Guía breve SRL, pp. 14-18 del PDF: márgenes, texto, títulos y aparato crítico.
% Criterios editoriales: docs/template-requirements.md.
\geometry{margin=3cm,headheight=25pt,headsep=8mm,footskip=12mm}
\setlength{\parindent}{1.25cm}
\setlength{\JustifyingParindent}{1.25cm}
\setlength{\parskip}{0pt}
\AtBeginDocument{\doublespacing\justifying}
\setcounter{secnumdepth}{3}
\setcounter{tocdepth}{3}
\renewcommand{\thesection}{\Roman{section}}
\renewcommand{\thesubsection}{\Alph{subsection}}
\renewcommand{\thesubsubsection}{\arabic{subsubsection}}
% Destinos PDF únicos aunque las letras y cifras se reinicien por sección.
\renewcommand{\theHsection}{\arabic{section}}
\renewcommand{\theHsubsection}{\theHsection.\arabic{subsection}}
\renewcommand{\theHsubsubsection}{\theHsubsection.\arabic{subsubsection}}

\titleformat{\section}[block]
  {\normalfont\normalsize\bfseries\centering\color{SRLHeadingColor}}
  {\thesection.}{0.5em}{}
\titleformat{\subsection}[block]
  {\normalfont\normalsize\bfseries\raggedright\color{SRLHeadingColor}}
  {\thesubsection.}{0.5em}{}
\titleformat{\subsubsection}[runin]
  {\normalfont\normalsize\bfseries\color{SRLHeadingColor}}
  {\thesubsubsection.}{0.5em}{}[.]
\titlespacing{\section}{0pt}{\baselineskip}{0.5\baselineskip}
\titlespacing{\subsection}{0pt}{0.5\baselineskip}{0pt}
\titlespacing{\subsubsection}{1.25cm}{0.5\baselineskip}{0.5em}

% Bosquejo sincronizado con los títulos del desarrollo, sin páginas ni puntos guía.
% Entradas del bosquejo sin negrita.
\titlecontents{section}[0pt]{\normalfont\normalsize}
  {\thecontentslabel.\enspace}{}{}
\titlecontents{subsection}[1.25cm]{\normalfont\normalsize}
  {\thecontentslabel.\enspace}{}{}
\titlecontents{subsubsection}[2.5cm]{\normalfont\normalsize}
  {\thecontentslabel.\enspace}{}{}

\pagestyle{fancy}
\fancyhf{}
\fancyhead[L]{\IfFileExists{\SRLHeaderLogo}{%
  \includegraphics[height=7mm,width=25mm,keepaspectratio]{\SRLHeaderLogo}}{}}
\fancyhead[R]{{\normalfont\fontsize{10}{12}\selectfont\thepage}}
\renewcommand{\headrulewidth}{0pt}
\renewcommand{\footrulewidth}{0pt}

% El estilo por defecto de Turabian se adapta a notas SRL de 10 pt,
% alineadas a la izquierda y con la misma familia del cuerpo.
\makeatletter
\AtBeginDocument{%
  \renewcommand{\footnotesize}{\@setfontsize\footnotesize{10}{12}}%
  \long\def\@makefntext#1{%
    \normalfont\footnotesize\setstretch{1}\RaggedRight
    \noindent\@makefnmark\hspace{0.4em}#1}%
}
\makeatother

% Cita larga: sin comillas/negrita, 12 pt, espacio sencillo, solo margen izquierdo.
\newenvironment{SRLblockquote}{%
  \par\begingroup\normalfont\normalsize\singlespacing
  \begin{list}{}{%
    \setlength{\leftmargin}{1cm}\setlength{\rightmargin}{0pt}%
    \setlength{\listparindent}{0pt}\setlength{\itemindent}{0pt}%
    \setlength{\parsep}{0pt}\setlength{\topsep}{0.5\baselineskip}}%
  \item[]\ignorespaces
}{\end{list}\endgroup\ignorespacesafterend}

% Portada: 12 pt, centrada, mayúsculas; logos opcionales.
\newcommand{\SRLmaketitle}{%
  \begin{titlepage}
    \thispagestyle{empty}\centering\normalfont\normalsize\doublespacing
    \IfFileExists{\SRLCoverLogo}{%
      \includegraphics[width=7cm,height=2.2cm,keepaspectratio]{\SRLCoverLogo}\par
    }{\vspace*{1.2cm}}
    \vfill
    {\color{SRLHeadingColor}\bfseries\MakeUppercase{\EssayTitle}\par}
    \ifdefempty{\EssaySubtitle}{}{%
      {\color{SRLHeadingColor}\MakeUppercase{\EssaySubtitle}\par}}
    \vfill
    \MakeUppercase{Estudiante: \EssayAuthor}\par
    \medskip
    \MakeUppercase{Docente: \EssayProfessor}\par
    \vfill
    {\color{SRLCoverDetailColor}\MakeUppercase{\EssayCourse}\par}
    \medskip
    \MakeUppercase{\EssayInstitution}\par
    \vfill
    \MakeUppercase{\EssayCity}\par
    {\color{SRLCoverDetailColor}\MakeUppercase{\EssayDate}\par}
  \end{titlepage}
  % Chicago/Turabian cuenta la portada aunque oculte su número.
  % titlepage reinicia el contador en documentos oneside; continuar en página 2.
  \setcounter{page}{2}
  \pagestyle{fancy}
}

% Separa las partes que la guía denomina «hoja única» y avisa si desbordan.
% El desbordamiento genera una advertencia para la revisión editorial.
\newenvironment{SRLsinglepage}[1]{%
  \clearpage\def\SRLPagePart{#1}\label{srl:#1:start}%
}{%
  \par\label{srl:\SRLPagePart:end}%
  \ifnum\getpagerefnumber{srl:\SRLPagePart:end}>\getpagerefnumber{srl:\SRLPagePart:start}\relax
    \PackageWarningNoLine{srl-template}{La parte '\SRLPagePart' excede una pagina.
      Revisar el requisito SRL de hoja unica}%
  \fi
  \clearpage
}
\makeatletter
\newcommand{\SRLmakeoutline}{%
  \begin{SRLsinglepage}{bosquejo}
    \section*{Bosquejo}
    \begingroup\normalfont\normalsize\color{black}\doublespacing
      \@starttoc{toc}%
    \endgroup
  \end{SRLsinglepage}
}
\makeatother
\defbibheading{srlbibliography}{\section*{Bibliografía}}
\renewcommand{\bibfont}{\normalfont\normalsize\doublespacing}
\setlength{\bibhang}{1.25cm}
\setlength{\bibitemsep}{0pt}
\setlength{\bibparsep}{0pt}

% Diagnóstico reproducible para CI y revisión del PDF.
\AtBeginDocument{%
  \typeout{SRL-LAYOUT: textwidth=\the\textwidth; textheight=\the\textheight; indent=\the\parindent}%
  \typeout{SRL-SPACING: baseline=\the\baselineskip}%
}

